Validierung
===========
-> was soll rauskommen, was soll das bewirken? (mögliche Antworten und was bedeuten diese Antworten?, ja nein und graustufen, wie geht man damit um?)

Interviews sind keine leichte Arbeit
-> müssen nicht gigantisch groß sein (viertel stunde Interview oder Umfrage)
-> am beispiel meines Teams

- wurden die User besser supportet durch die Funktion?

Wie funktionieren Interviews
wie sorge ich dafür dass ein Interview wissenschaftlich korrekt durchgeführt wird?
-> in die BA: Interview wurde nach Schema X nach Quelle Y durchgeführt
-> möglichst alle 15 zusammenbekommen
-> vorher ankündigen

ich hab eine Technologie, hab das hiermit Umgesetzt + Validierung durch ein Interview, wurden eure Ansprüche umgesetzt und die Anforderungen erfüllt (zur Validerung
-> Lernen: wie werden Interviews durchgeführt?
-> Interviewmethoden und Namen dazu

Umfrage / Fragebogen / Experteninterview - wo sie einen Nutzen und Probleme an dem Dashboard sehen
-> allgemein (potenziell vorher) - was haltet ihr von der Idee, würdetihr das Nutzen?
-> explizit auf das Dashboard (nach Vorstellung) - hier das ist meine Umsetzung, ist das so, wie ihr euch das vorgestellt habt?
- Feedback einbauen und nochmal auswerten?

-	Prof: „Interviews sind das heftigste was man für eine BA machen kann, jeder Prof wird immer ein Fehler in der Methodik finden“
o	Für jede Frage: ein festes Anfragenspektrum definiert wird
o	Alle möglichen Antworten müssen qualitativ eingeordnet werden können
o	Offene Fragen mit unerwarteten Antwortmöglichkeiten vermeiden
-> z.B. YME oder Sasi fragen, was man auf jede Frage antworten könnte

% ****************************************************************************
%           Abgrenzung Evaluation
% ****************************************************************************

Evaluation
==========
fällt unter Wirtschaftsinformatik
=================================
wie sieht  das alte Tool aus
welche Anforderungen wurden durch das neue Tool umgesetzt
-> quantitative und qualitative Maßnahmen
Unterschiede Herausarbeiten
-> These: sind die Anforderungen, die vorher existiert haben, umgesetzt wurden
-> Problem: wie möchte ich das messen

ist ne Validierung
-> warum gibt es das Dashboard, welches Problem wurde gelöst

% ****************************************************************************

- Fragen testen in anderem Kreis
- Feedback aus anderem Kreis einholen
- qualitativ und quantitativ kombinieren
-> nur im Bezug auf "gibt es Anmerkungen, willst du das weiter erläutern?"
https://de.surveymonkey.com/mp/survey-question-types/
- später auf narrative Interviews und grundsätzlich qualitative Interviews, Sozialwissenschaften, Experteninterview eingehen
https://link.springer.com/book/9783531172385


Mein Ziel - was soll rauskommen, was soll das bewirken?
•	wurden die User besser supportet durch die Funktion?
•	
Bewerten von Antworten
•	mögliche Antworten und was bedeuten diese Antworten?
-> ja nein und graustufen
-> wie geht man damit um?
Interviews sind keine leichte Arbeit
-> müssen nicht gigantisch groß sein (viertel stunde Interview oder Umfrage)
-> am Beispiel meines Teams
Wie funktionieren Interviews
•	wie sorge ich dafür dass ein Interview wissenschaftlich korrekt durchgeführt wird?
•	-> in die BA: Interview wurde nach Schema X nach Quelle Y durchgeführt
•	-> möglichst alle 15 zusammenbekommen
•	-> vorher ankündigen
•	intuitiv Interview
- 15 -4, an die musste ich reporten

offizielle Fragen:
vor dem zeigen
1.	Person Einordnen: wie oft schaust du nach den Testpool Ergebnissen
-> welche davon? (nur bestimmte, alle …)
a.	daily, basically daily
b.	regelmäßig
i.	oft (z.B. 1x die Woche)
ii.	selten (z.B: 1x im Monat)
iii.	wenn ich ne Änderung vorgenommen habe
c.	unregelmäßig
i.	wenn was schief läuft
ii.	once a while
d.	quasi nie / nicht meine Aufgabe
2.	allgemein: (potenziell vorher)
a.	was haltet ihr von der Idee, würdet ihr das Nutzen?
 vorsicht suggestiv Frage!!!
b.	was sind deine Ansprüche / Anforderungen an ein Dashboard?
Test durchführen in bisheriger Umgebung
Dashboard zeigen
3.	Standardisierter Test „Fehler gefunden“
-> wo sind Fehler
-> gib mir die Übersichtsseite von diesem System / Test && einmal mit Dashboard
a.	Ergebnisse einordnen – nach wie oft nutzt man die Übersicht
b.	Vergleich Ausgangslage:
i.	Browser Clean wie er im Alltag genutzt wird – mit Bookmarks und allem
ii.	Dashboard Link im Teams Chat
-> Zeit zum Öffnen der Webseite abziehen
-> Dashboard nicht bekannt / zu sehen
4.	Erster Eindruck: was fällt dir zuerst auf, was springt dir ins Auge, wo ist der Fokus
5.	allgemeine Einschätzung:
a.	wirst du das Dashboard anstelle der bisherigen Lösung verwenden?
-> wie beeinflusst das Dashboard das Monitoring?
 Overhead, Erleichterung, Neutral? – 5-stufig
b.	wird das Dashboard etwas daran ändern, wie häufig du in die Testpool-Ergebnisse schaust?
6.	Rückkopplung auf vorher: 
-> wurden deine Ansprüche/Anforderungen dMn umgesetzt / getroffen?
7.	Feedback: was ist dir noch jetzt aufgefallen im Umgang
a.	spezifische Sachen die gefallen … die nicht gefallen
b.	positiv, negativ
c.	was sind Nutzen und Probleme?

https://de.surveymonkey.com/mp/survey-question-types/ 
-> Lernen: wie werden Interviews durchgeführt?
-> Interviewmethoden und Namen dazu

Prof: „Interviews sind das heftigste was man für eine BA machen kann, jeder Prof wird immer ein Fehler in der Methodik finden“
•	Für jede Frage: ein festes Anfragenspektrum definiert wird
•	Alle möglichen Antworten müssen qualitativ eingeordnet werden können
•	Offene Fragen mit unerwarteten Antwortmöglichkeiten vermeiden
-> z.B. YME oder Sasi fragen, was man auf jede Frage antworten könnte

 
Interview - Leitfaden, strukturiert, narrativ, halb-strukturiert
Überhaupt wird häufig auch die theoretische und methodologische Begründung für die Auswahl qualitativer Verfahren "vergessen".
-> qualitativ ist zu umfangreich für die Arbeit
-> quantitativ it is
Vermeiden Sie Suggestivfragen
Formulieren Sie Fragen, die die Befragten nicht zu einer bestimmten Antwort verleiten oder beeinflussen. Anstatt zu fragen: „Wie gut war Ihre Erfahrung mit unserem Kundenservice?“, fragen Sie: „Wie würden Sie Ihre Erfahrung mit unserem Kundenservice bewerten?“

